\section{A further positive example requiring a nonsimplex extension}

\subsection{O-023: the order-ideal extension improves every prefix-simplex test}
Keep
$m=2^{200}$, $r=m^2$, $n=m^{25}$, and $d=3n$,
and now define
\[
 \lambda=(nd-26r-1000000+50,\ 1000000,\ 28,\ 26^{r-3}).
                                                               \tag{O-023.1}
\]
Then
\[
                 4\leq q^m_\lambda(d[n])
                               \leq K_n(\lambda)=a_\lambda(d[n]).
                                                               \tag{O-023.2}
\]
Every choice of parameters in the prefix-simplex criterion
M-017 fails to reach this shape at size $n$.
The exact support-cover and one-cut failures of O-022
also remain. Thus M-018 makes a further strict reduction
of the proved candidate conditions.
No exact numerical value of the multiplicities in
(O-023.2) is asserted.

\paragraph{The nonsimplex extension fits.}
For the order ideal
\[
 A=\{x_2^a x_3^c:\ a+c\leq999974,\quad c\leq28\}
\]
one has
\[
       |A|=29\cdot999975-\frac{28\cdot29}{2}
                          =28998869.
                                                               \tag{O-023.3}
\]
The full nested support of (O-023.1) has total degree
at most one million, suffix degree from coordinate three
at most twenty-eight, and suffix degree from coordinate
four at most twenty-six.
If its monomial degree exceeds twenty-six, choose an
early submonomial in $x_2,x_3$ of the excess degree.
There are enough early factors because the later suffix
has degree at most twenty-six. Its degree is at most
$999974$, and its $x_3$ exponent is at most twenty-eight,
so it lies in $A$. The remaining monomial has degree
twenty-six. Smaller monomials use the zero element of $A$.
Thus the full support is contained in $A+D_{26}$.

Use M-018 with $b=26$ and
$T=m\binom{m^2+11}{12}$.
The size is $28998870T<n$, because
$28998870<C_*=478126755$ and the uniform normalized
estimate of M-017 applies. This proves $K_n=a_\lambda$.

\paragraph{Four positive permanent copies.}
Multiply the four degree-six binary seed classes of M-012
by
\[
             W_r^{13}\,W_3\,W_2^{499983}\,
                         c_0^{\,d-13r-999975}.          \tag{O-023.4}
\]
The exponent is nonnegative. O-011 supplies one basis
at a permanent point where every displayed leading
Hessian minor and $c_0$ is nonzero. Therefore the common
factor is nonzero in the integral coordinate algebra of
the irreducible permanent closure, and multiplication
preserves the four independent seed classes.
Their degree is
$6+13r+3+2\cdot499983+(d-13r-999975)=d$.
Their second weight is
$6+26+2+999966=1000000$, their third weight is $28$,
and all later weights are $26$. The first weight is
the remaining part $nd-26r-999950$, as required.
This proves $q^m_\lambda\geq4$ and hence (O-023.2).

\paragraph{All prefix-simplex choices fail.}
Consider any M-017 parameter choice that could cover this
highest-weight support. Its even base order is $b=2\ell+2$,
and it must satisfy $\lambda_2\leq b+t$ and
$\lambda_k\leq b$. Taking more ambient variables than
the actual length $r$ only increases its size; it suffices
to give lower bounds using this length.

If $b\geq28$, then $\ell\geq13$, and even the unextended
core has size at least
\[
     2m\binom{r+12}{13}
           \geq\frac{2m^{27}}{13!}>m^{25}=n.             \tag{O-023.5}
\]
If $b=26$, the first possible cut is $k=4$, since
$\lambda_3=28$. Therefore $p=k-2\geq2$ and
$t\geq999974$. Its size is at least
\[
 \frac{m^{25}}{12!}
       \left(1+\binom{999976}{2}\right)>m^{25},
                                                               \tag{O-023.6}
\]
because $T\geq m^{25}/12!$ and the parenthesized integer
is larger than $12!$.
Finally, if $b\leq24$, every actual part is larger than
$b$, so the cut must be at or after $r+1$.
Then $p\geq r-1$, $t\geq13$, and $T\geq m$.
Its size is at least
\[
                    m\binom{r+12}{13}
                          \geq\frac{m^{27}}{13!}>n.    \tag{O-023.7}
\]
These cases exhaust all allowed even base orders.
The strict inequalities use $m^2>13!$, and are verified
by the exact certificate. Thus no M-017 parameter choice
satisfies its size condition.

\paragraph{The earlier failures persist.}
The nested support (O-022.6) is contained in the present
support: its total degree is at most $32$, its suffix
from the third coordinate is at most $26$, and therefore
its suffix from the fourth coordinate is at most $26$.
In particular the present support contains all degree-$26$
forms in the last $r-2$ variables. The exact-complexity
argument (O-022.7) again forces every F-018 exact support
cover to have size greater than $n$.

Each part of (O-023.1), except the first, is at least
the corresponding part of (O-022.1).
Every one-cut family and coordinate box is consequently
an enlargement of its O-022 counterpart. The border-valid
dimension failures (O-022.8)--(O-022.10) and the box
size failure (O-022.11) remain valid.
The full point-jet target is still too large, since it
contains the degree-$27$ monomials used there.
The shape meets the BIP staircase, A-025's joint gate,
and the M-017 coupled necessary conditions; the complete
M-017 size failure was proved above.
Its tail is $26r+999950\leq md$ and its length is $m^2$.
The shape-specific Chow threshold is at least
$d-1=3n-1>n$; the uniform threshold and Grenet failure
are unchanged. O-020 does not apply because $\lambda_3=28$.

These checks concern exact proved exclusions, including
ones with actual multiplicity at least four. They do not
claim that (O-023.1) is an obstruction: (O-023.2)
proves that it is not.

\paragraph{Geometric dictionary.}
The set $A$ is a finite monomial order ideal in an integral
torus-graded coefficient module, as in SGA~3 I,
Proposition~4.7.3, pp.~35--36. M-018 gives its actual
universally injective coordinate pullback. The full
highest-weight support is the Schur/exterior calculation
of F-016, with SGA~6 VI, Theorem~3.3, pp.~392--393,
and Macdonald I(4.3), p.~63, Appendix~A(5.3)--(5.6),
pp.~155--156. The domain multiplication retains four
sections of the actual orbit-closure coordinate algebra.
The exact-complexity obstruction and the separate
border dimension arguments retain the distinct scopes
spelled out in O-022. No saturation statement follows
from a virtual class alone.
